\documentclass[11pt,a4paper]{article}
\usepackage[T1]{fontenc}
\usepackage[utf8]{inputenc}
\usepackage[main=italian,provide=*]{babel}
\usepackage{amsmath,amssymb}
\usepackage{geometry}
\geometry{margin=2.5cm}
\usepackage{booktabs}
\usepackage{seqsplit}
\usepackage{hyperref}
\hypersetup{colorlinks=true,linkcolor=blue,citecolor=blue,urlcolor=blue}

\title{Manuale d'uso --- correlazioni dinamiche trimero anello\\
\large versione attuale (circuito VQE reale \texttt{W-2q.6})}
\author{Note di lavoro --- Tesi triennale, Samuele Schiavi\\ Relatore: Prof.\ Alessandro Chiesa}
\date{}

\begin{document}
\maketitle

\section*{Scopo}

Questo manuale spiega come eseguire la simulazione delle correlazioni dinamiche
$C_{ij}^{\alpha\beta}(t)$ per il trimero ad anello nella versione
\emph{attuale}: lo stato fondamentale non è più preparato con le ampiezze
esatte, ma con il circuito VQE reale (ansatz \texttt{W-2q.6}), ottimizzato
numericamente. È la versione da usare per tutto ciò che deve rappresentare
fedelmente "quello che gira su un dispositivo/simulatore quantistico", dato
che \texttt{prepare\_state} non è un'operazione fisicamente implementabile
in generale. La versione a stato esatto resta disponibile (è ancora il
comportamento di default se non si passano parametri VQE) e serve da
riferimento: vedi il manuale gemello
\texttt{\seqsplit{manuale\_uso\_correlazioni\_trimero\_anello\_statevector.tex}}.

\section*{Prerequisiti}

Tutti i file del manuale gemello, più:
\begin{itemize}
\item \texttt{vqe\_w2q6\_trimero\_anello.py} --- costruzione dell'ansatz
\texttt{W-2q.6} e ottimizzazione VQE.
\item \texttt{validate\_vqe\_circuito\_correlazioni.py} --- confronto sistematico
contro la versione a stato esatto.
\item \texttt{scan81\_vqe\_trimero\_anello.py} --- scan completo delle 81
combinazioni (statevector + shot), preparazione VQE (Passo 4).
\end{itemize}
Stesso punto di lavoro di riferimento: $J=1$, $J'=0.4$, $b=b_c=2.4$,
$D=0.15$ (Opzione B), $t=1.3$, $N=200$ passi di Trotter.

\begin{quote}
\textbf{Importante:} i parametri ottimali dell'ansatz dipendono dal punto di
lavoro $(J,J',b,D)$. Se si cambia uno di questi valori bisogna
\emph{rieseguire} l'ottimizzazione (Passo 1) prima di usare
\texttt{ansatz\_params} altrove --- non riusare parametri ottenuti a un
punto diverso, la fedeltà non è garantita.
\end{quote}

\section*{Passo 1 --- ottimizzazione VQE (una tantum per punto di lavoro)}

\begin{verbatim}
python vqe_w2q6_trimero_anello.py
\end{verbatim}
Esegue 12 partenze casuali (COBYLA, seed fissato) più un polish finale
(L-BFGS-B), e stampa l'energia trovata a ogni partenza. Al punto di lavoro
standard tutte convergono allo stesso minimo:
\begin{align*}
E_{\rm vqe} &= -5.53437001739340, \qquad E_{\rm esatto} = -5.53437001739350,\\
\Delta E &= 9.8\times10^{-14}, \qquad \mathcal F = 0.99999999999961.
\end{align*}
Se il tuo output mostra $\Delta E$ o $1-\mathcal F$ molto più grandi (es.
$10^{-3}$ o peggio), l'ottimizzatore è rimasto bloccato in un minimo locale:
aumenta \texttt{N\_START} nello script o cambia il seed
(\texttt{np.random.default\_rng(0)}). Salva il risultato in
\texttt{w2q6\_params\_optimal.npz} (6 parametri reali, in radianti).

\section*{Passo 2 --- un correlatore singolo, con preparazione VQE}

\begin{verbatim}
import numpy as np
from circuito_correlazioni_trimero_anello import correlator_from_circuit

data = np.load("w2q6_params_optimal.npz")
params = data["params"]

J, Jp, b, D = 1.0, 0.4, 2.4, 0.15
C = correlator_from_circuit(1, 'y', 1, 'y', t=1.3, N=200, J=J, Jp=Jp, b=b, D=D,
                             ansatz_params=params)
print(C)
\end{verbatim}
Passando \texttt{ansatz\_params}, la preparazione dello stato usa il circuito
\texttt{w2q6\_circuit()} (X + tre blocchi \texttt{W} + $R_y$ finali) invece di
\texttt{prepare\_state}. \texttt{psi0} non serve più in questo caso (viene
ignorato se passato insieme a \texttt{ansatz\_params}).

\section*{Passo 3 --- validazione contro la versione a stato esatto}

\begin{verbatim}
python validate_vqe_circuito_correlazioni.py
\end{verbatim}
Richiede che \texttt{w2q6\_params\_optimal.npz} sia già presente (Passo 1).
Ricalcola tutte e 81 le combinazioni $C_{ij}^{\alpha\beta}(t)$ sia con
preparazione VQE sia con preparazione esatta, e stampa l'errore medio e
massimo tra le due. Valori attesi al punto di lavoro standard: errore medio
$3.0\times10^{-7}$, massimo $8.3\times10^{-7}$ (su $C_{21}^{zz}$) --
dell'ordine di $\sqrt{1-\mathcal F}$, come atteso per un residuo di fedeltà
a livello di stato propagato agli osservabili dinamici. Errori molto più
grandi indicano parametri VQE non ottimizzati bene (torna al Passo 1).

\section*{Passo 4 --- scan completo delle 81 combinazioni (statevector + shot)}

\begin{verbatim}
python scan81_vqe_trimero_anello.py
\end{verbatim}
Mirror di \texttt{scan81\_trimero\_anello.py} (manuale gemello, Passo 4), con
la stessa struttura ma preparazione VQE al posto di \texttt{prepare\_state}.
Richiede \texttt{w2q6\_params\_optimal.npz} già presente (Passo 1). Stampa lo
stesso tipo di riepilogo (zeri strutturali, errori medi/massimi) e salva
\texttt{scan81\_vqe\_results.npz}/\texttt{.json} nella cartella corrente ---
percorsi relativi, quindi va eseguito da dentro quella cartella. Valori
attesi al punto di lavoro standard: errore totale (Trotter+VQE) medio
$8.81\times10^{-4}$, massimo $2.64\times10^{-3}$ --- indistinguibile,
a questa precisione, da quello della versione a stato esatto; errore
statistico (shot vs statevector) medio $1.33\times10^{-2}$, massimo
$3.33\times10^{-2}$, ancora dominante sull'errore di Trotter e sul residuo
VQE. Serve come dato precomputato alla Sez.~7 di
\texttt{\seqsplit{circuito\_correlazioni\_trimero\_anello\_vqe.ipynb}} (vedi
"Come esplorare tutto insieme" sotto).

\section*{Come esplorare tutto insieme: il notebook}

\texttt{\seqsplit{circuito\_correlazioni\_trimero\_anello\_vqe.ipynb}} esegue
in sequenza tutti i passi sopra (ottimizzazione inclusa, generata inline,
non serve rieseguire il Passo 1 a parte) e in più mostra, con la
preparazione VQE al posto di \texttt{prepare\_state}: heatmap $|C|$
d'insieme, vista nel tempo su griglia $9\times9$, selettore per ispezione
singola, verifica della simmetria $C_{11}=\eta_\alpha\eta_\beta\,C_{22}$,
spettro dietro i casi ricco/piatto, confronto con lo scan a shot finiti
(Passo 4, va eseguito prima per generare \texttt{scan81\_vqe\_results.npz}).
È il modo più diretto per farsi un'idea complessiva senza lanciare gli
script uno per uno; gli script restano il modo giusto per rieseguire un
singolo pezzo o integrare il codice altrove.

\section*{Cosa NON cambia rispetto alla versione a stato esatto}

\texttt{build\_correlator\_circuit} e \texttt{correlator\_from\_circuit} sono
retrocompatibili: se non si passa \texttt{ansatz\_params} (o si passa
\texttt{None}, il default), il comportamento è identico, bit per bit, alla
versione statevector --- tutto il codice esistente
(\texttt{scan81\_trimero\_anello.py}, \texttt{\seqsplit{validate\_shot\_noise\_trimero\_anello.py}},
e i primi tre notebook della fase, che usano ancora \texttt{prepare\_state})
continua a funzionare senza modifiche.

\section*{Problemi comuni}

\begin{itemize}
\item \texttt{FileNotFoundError} su \texttt{w2q6\_params\_optimal.npz} ---
va generato eseguendo \texttt{vqe\_w2q6\_trimero\_anello.py} (Passo 1); non è
incluso di default (dato binario specifico del punto di lavoro, non versionato
nel Project).
\item Fidelity bassa dopo l'ottimizzazione --- vedi nota al Passo 1
(minimo locale); aumentare le partenze multistart di solito risolve.
\item Import circolare se si prova a importare \texttt{ground\_state} da
\texttt{circuito\_correlazioni\_trimero\_anello.py} dentro
\texttt{vqe\_w2q6\_trimero\_anello.py} --- per questo il secondo modulo
definisce una propria copia locale \texttt{\_ground\_state} (stessa logica,
duplicata apposta per evitare l'import circolare tra i due file).
\item \texttt{FileNotFoundError} su \texttt{scan81\_vqe\_results.npz} nella
Sez.~7 del notebook --- va prima generato eseguendo
\texttt{\seqsplit{scan81\_vqe\_trimero\_anello.py}} (Passo 4), che a sua volta
richiede \texttt{w2q6\_params\_optimal.npz} (Passo 1).
\end{itemize}

\section*{Riferimenti}

Derivazione e giustificazione dell'ansatz: notebook
\texttt{\seqsplit{confronto\_ansatz\_entangler\_trimero\_anello.ipynb}},
\texttt{\seqsplit{analisi\_espressivita\_PMA\_anello.ipynb}}, e
\texttt{\seqsplit{trimero\_anello\_vqe\_dm.tex}}. Documentazione teorica
aggiornata: \texttt{\seqsplit{circuito\_correlazioni\_trimero\_anello\_spiegato.tex}}.
Catalogo completo della fase: \texttt{\seqsplit{trimero\_anello\_correlazioni.tex}}.
Cronologia: \texttt{log\_decisioni.md}, sezione "VQE reale nel circuito dei
correlatori (trimero anello)".

\end{document}